\section{Related Work}
\label{sec:related_work}

Our work builds on and connects three main streams of literature: parameter-space model merging, geometric/manifold-based model alignment, and test-time adaptive optimization.

\subsection{Linear Parameter Merging}
Model merging seeks to combine separate specialized models into a single multi-task model without training from scratch~\cite{choshen2022fusing, donyehiya2023cold}. 
Early formulations focused on simple parameter-wise linear combinations. 
Model Soups~\cite{wortsman2022model} averages weights of fine-tuned models trained on the same task to improve generalization. 
Task Arithmetic~\cite{ilharco2022editing} computes task-specific updates (task vectors) and adds them linearly to the pre-trained base model. 
To handle parameter interference, RegMean~\cite{jin2023regmean} optimizes merging weights via feature-correlation matrices, and Fisher Merging~\cite{matena2022merging} weights parameters by their diagonal Fisher information. 
More recently, Ties-Merging~\cite{yadav2023resolving} addresses sign conflicts by truncating small values, finding sign consensus, and performing elected averaging. 
DARE~\cite{yu2023language} uses extreme sparsification to mitigate interference. 
While effective, these methods are restricted to Euclidean space and suffer from severe performance degradation when task-specific parameters lie in disjoint basins separated by high-loss ridges.

\subsection{Geometric and Manifold Model Merging}
To overcome the limitations of flat linear interpolation, researchers have explored the non-linear structure of parameter spaces. 
One prominent line of work studies linear mode connectivity (LMC) modulo permutation symmetries~\cite{entezari2021role, garipov2018loss}. 
Git Re-Basin~\cite{ainsworth2022git} and related alignment methods~\cite{jordan2022repair, juneja2023linear} leverage Hungarian algorithms or gradient descent to find permutation matrices that align different neural networks into a shared basin. 
ZipIt!~\cite{zhang2023zipit} merges models with disparate features by matching and combining redundant neurons. 
While powerful, these methods are limited to discrete permutation matrices and do not generalize to continuous deformations. 
OrthoMerge~\cite{yang2026orthomerge} addresses geometric preservation by mapping parameter updates to the Lie algebra $so(d)$ of the orthogonal group. 
However, OrthoMerge relies on a predefined, rigid Riemannian manifold projection. 
In contrast, FoldMerge (Neural Origami) is the first to employ highly expressive, learned, and data-driven coordinate transformations using continuous weight-space diffeomorphisms.

\subsection{Test-Time Adaptive Merging}
Recent methods optimize merging parameters on unlabeled downstream validation streams. 
AdaMerging~\cite{yang2024adamerging} adaptively optimizes task-wise coefficients at test-time by minimizing entropy. 
Representation Surgery~\cite{gu2024patch} performs feature-space projection alignment to minimize inter-task interference. 
SyMerge~\cite{jung2025symerge} introduces an unsupervised self-labeling test-time adaptation protocol where task-specific parameters are adaptively scaled via expert teacher guidance. 
Although SyMerge achieves strong multi-task performance, its adaptation is still fundamentally limited to linear scaling of the classifier heads. 
FoldMerge builds upon the self-labeling protocol of SyMerge, but introduces a radically different, non-linear coordinate warping framework. 
Rather than performing linear scale-tuning, we warp the entire weight-space coordinate system via a multi-layer normalizing flow, yielding significantly richer representational capacity.
